%% ============================================================================
%% Fichier  : annexes/Ann_B_outils.tex
%% Titre    : Référentiel des outils
%% Version  : 1.0 -- novembre 2026
%% ============================================================================
%% Alignée sur les annexes des autres cours de l'auteur (ACSI, INF3226).
%% ============================================================================

\chapter{Référentiel des Outils}
\label{ann:outils}

Cette annexe rassemble les outils mobilisés dans le cours et dans les labs de
NDAP SUKULU. Les noms de logiciels sont des exemples~: les versions évoluent, et
tout rapport cite l'outil \emph{et sa version}. Un outil se valide avant
usage, sur des données de référence (ISO/IEC~27041).

\section*{Systèmes et mémoire}

\begin{tableaucours}{L{3.6cm}L{6.2cm}L{2.6cm}}{%
    \tableauHeader{Outil} &
    \tableauHeader{Usage} &
    \tableauHeader{Chapitre}
}
Outils d'analyse du registre et de NTFS (\$MFT, \$UsnJrnl) & Preuves d'exécution, chronologie des fichiers & Ch.~\ref{chap:for-sys} \\
Plaso (\texttt{log2timeline}) & Timeline unique de toutes les sources & Ch.~\ref{chap:for-sys} \\
Volatility~3 & Processus, connexions, injections en mémoire & Ch.~\ref{chap:for-sys} \\
\end{tableaucours}

\section*{Réseau, mobile, nuage, objets}

\begin{tableaucours}{L{3.6cm}L{6.2cm}L{2.6cm}}{%
    \tableauHeader{Outil} &
    \tableauHeader{Usage} &
    \tableauHeader{Chapitre}
}
Wireshark, Zeek & Flux, sessions, beaconing, exfiltration & Ch.~\ref{chap:for-res} \\
Lecteurs SQLite, Autopsy & Messages, appels, Mobile Money, EXIF & Ch.~\ref{chap:for-mob} \\
Journaux du fournisseur (par exemple CloudTrail) & Appels d'interface, fenêtres aveugles & Ch.~\ref{chap:for-cloud} \\
\texttt{binwalk} & Extraire et lire un micrologiciel & Ch.~\ref{chap:for-iot} \\
\texttt{kubectl}, journaux d'audit & Conteneurs et grappes Kubernetes & Ch.~\ref{chap:containers} \\
Outils OSINT (WHOIS, archives du web, explorateurs de chaînes) & Domaines, profils, cryptoactifs & Ch.~\ref{chap:for-osint} \\
\end{tableaucours}

\section*{Cryptographie et méthodes formelles}

\begin{tableaucours}{L{3.6cm}L{6.2cm}L{2.6cm}}{%
    \tableauHeader{Outil} &
    \tableauHeader{Usage} &
    \tableauHeader{Chapitre}
}
Bibliothèques post-quantiques (Open Quantum Safe) & Expérimenter ML-KEM et ML-DSA & Ch.~\ref{chap:quantique} \\
Vérificateurs de protocoles (Tamarin, ProVerif) & Analyser formellement un protocole & Ch.~\ref{chap:form-meth} \\
\end{tableaucours}

\begin{boxalert}
\textbf{Rappel.} Aucun outil n'est utilisé sur un système réel sans cadre
légal ni autorisation écrite (loi n°~2010/012, art.~65, 69 et~86~;
Annexe~\ref{ann:lois}).
\end{boxalert}
